\section{Введение}

В условиях современного развития информационных технологий и усложнения программных продуктов существенно возрастает сложность процессов их разработки и сопровождения. Современные компании, особенно в сфере информационных технологий, всё чаще переходят от работы отдельных автономных команд к модели, в которой несколько команд параллельно развивают один продукт или платформу. Такая организация труда требует более высокого уровня координации, прозрачности процессов и согласованности действий между участниками.

В этих условиях особое значение приобретают гибкие методологии управления проектами, ориентированные не только на локальную эффективность отдельных команд, но и на оптимизацию всей системы в целом. Одной из наиболее распространённых и практически применяемых методологий является Kanban. В отличие от итеративных подходов, таких как Scrum, Kanban ориентирован на непрерывный поток работ, визуализацию процессов и ограничение количества незавершённых задач. Данный подход хорошо масштабируется и широко используется в организациях со сложной структурой и уже существующими продуктами.

При этом ключевые принципы Kanban, включая ограничение WIP (Work in Progress), управление потоком и анализ метрик производительности, достаточно сложно усваиваются при изучении исключительно теоретических материалов. Без практического опыта участникам трудно осознать влияние управленческих решений на пропускную способность системы, время выполнения задач и стабильность процессов. В результате между формальным знанием методологии и её реальным применением на практике возникает существенный разрыв.

Для преодоления данного разрыва в корпоративном обучении и образовательных программах широко применяются бизнес-игры и игровые симуляции. Они позволяют участникам в безопасной среде экспериментировать с процессами, принимать управленческие решения и наблюдать их последствия без риска для реальных проектов. Игровой формат способствует более глубокому пониманию системных эффектов, возникающих при управлении потоком работ, и позволяет наглядно продемонстрировать влияние ограничений и изменений в процессе разработки.

Одной из таких симуляций является игра Featureban, моделирующая процесс разработки продукта с использованием Kanban-подхода. Классическая версия игры ориентирована на работу одной команды и позволяет продемонстрировать базовые принципы управления потоком. Мультикомандная версия Featureban представляет собой развитие данной симуляции и направлена на моделирование взаимодействия нескольких команд, одновременно работающих над общим продуктом. В рамках такой игры проявляются межкомандные зависимости, узкие места системы, проблемы синхронизации и эффекты локальной оптимизации, характерные для реальных крупных проектов.

Актуальность данной работы обусловлена возрастающей сложностью управления процессами разработки в условиях масштабирования продуктов и увеличения количества команд, вовлечённых в совместную работу. На практике многие организации сталкиваются с тем, что методы и практики Kanban, успешно применяемые на уровне одной команды, теряют эффективность при переходе к мультикомандному формату. При этом существующие инструменты и обучающие материалы не позволяют в полной мере продемонстрировать системные эффекты, возникающие при межкомандном взаимодействии, и обеспечить корректное воспроизведение игровой механики мультикомандной симуляции.

В этой связи актуальной является задача создания специализированного программного инструмента, позволяющего наглядно моделировать мультикомандные процессы разработки, демонстрировать влияние управленческих решений на всю систему в целом и обеспечивать автоматизацию игровых правил и сбора метрик. Решение данной задачи способствует повышению эффективности обучения и более осознанному внедрению Kanban-подходов в организациях со сложной структурой.

Целью данной курсовой работы является разработка онлайн-версии бизнес-игры «Мультикомандный Featureban», предназначенной для моделирования межкомандного взаимодействия и использования в обучающих и тренинговых форматах. Для достижения поставленной цели в рамках работы проводится анализ предметной области, формулируются требования к разрабатываемой системе, проектируется архитектура веб приложения и реализуются основные компоненты системы.

В результате планируется создать веб приложение, поддерживающее проведение игровых сессий с участием нескольких команд, разграничение ролей пользователей, визуализацию процессов и автоматизированный сбор ключевых метрик игрового процесса. Новизна данной работы заключается в разработке специализированного инструмента для онлайн-проведения мультикомандной версии Featureban, который учитывает особенности игровой механики, поддерживает сценарии симуляции и обеспечивает централизованное управление игровым процессом. Практическая значимость работы состоит в возможности использования разработанного приложения для проведения корпоративных тренингов и образовательных мероприятий, направленных на изучение принципов управления потоком и межкомандного взаимодействия.